% !PRINTED 94
\weberpage{94}
Then, as one recognises from
\[
\frac{\omega'_2}{\omega'_1} = \frac{c\omega_1+\partial\omega_2}{a\omega_1+b\omega_2}
\]
by separation of the real from the imaginary part, the
imaginary part of $\omega_2:\omega_1$ has the same sign as that of
$\omega'_2:\omega'_1$ (the positive one).

If we set
\weq{7}{\begin{gathered}
a_1 = \partial a'_1 - b a'_2 \\
a_2 = -c a'_1 + a a'_2
\end{gathered}}
\weq{8}{\begin{gathered}
b_1 = \partial b'_1 - b b'_2 - m'b\partial \\
b_2 = -cb'_1 + ab'_2 - m'ac,
\end{gathered}}
then one concludes from (3) that the function $T(u)$ satisfies not only
the conditions (1), but also the equations that arise from (1) by the
interchange of $\omega'_1$, $\omega'_2$, $a'_1$, $a'_2$, $b'_1$, $b'_2$ with $\omega_1$, $\omega_2$, $a_1$, $a_2$, $b_1$, $b_2$,
that is, the equations (1), § 17. It is
therefore at the same time a $T$-function of the periods $\omega'_1$, $\omega'_2$ and of the
periods $\omega_1$, $\omega_2$, which we indicate by the following equation:
\weq{9}{T'(u,\,\omega'_1\,\omega'_2) = T(u,\,\omega_1,\,\omega_2).}
% [sic] German source omits the comma between the two primed periods.
But by (5) and (7) we have
\[
a_2\omega_1 - a_1\omega_2 = (a'_2\omega'_1 - a'_1\omega'_2)(a\partial - bc),
\]
hence, if $m$ is the order of $T$,
\weq{10}{m = m'n.}
By (8), if $(g'_1, g'_2)$ is the characteristic of $T'$ (§ 17), the
characteristic $(g_1, g_2)$ of $T$ is
\weq{11}{(g_1, g_2) = (\partial g'_1 - b g'_2 - m'b\partial,\ -c g'_1 + a g'_2 - m'ac).}
\gesperrt{By the transformation of the $T$-functions one understands
the representation of the functions $T'$ with the periods
$\omega'_1$, $\omega'_2$ by $T$-functions with the periods $\omega_1$, $\omega_2$.}

The numbers $a$, $b$, $c$, $\partial$ are called the \gesperrt{transformation numbers}
and $n = a\partial - bc$ the \gesperrt{degree of transformation}.

In order to survey the form of this representation more clearly,
we shall seek the conditions under which $T'(u, \omega_1, \omega_2)$
becomes a $\Theta$-function of the $m$th order $\Theta(u,\omega)$ (§ 20). We
take $b'_1$, $b'_2$, and consequently also $b_1$, $b_2$, to be integers, so that
$b_1$, $b_2$, $b'_1$, $b'_2$ can be replaced by $g_1$, $g_2$, $g'_1$, $g'_2$. We
then have
\[
\omega_1 = 1, \quad \omega_2 = \omega, \quad a_1 = 0, \quad a_2 = m
\]
%CONT>>
